\documentclass[12pt]{article}

% ─── Packages ──────────────────────────────────────────────────────────────────
\usepackage[margin=1in]{geometry}
\usepackage{amsmath, amssymb, amsthm}
% \usepackage{physics}  % not available; using manual macros
\usepackage{booktabs}
\usepackage{array}
\usepackage{xcolor}
\usepackage{hyperref}
\usepackage{cleveref}
\usepackage{parskip}          % paragraph spacing instead of indent
% \usepackage{microtype}  % disabled for compatibility
\usepackage[font=small,labelfont=bf]{caption}

% ─── Theorem-like environments ─────────────────────────────────────────────────
\theoremstyle{definition}
\newtheorem{assumption}{Assumption}
\newtheorem{definition}{Definition}
\newtheorem{result}{Result}

% ─── Convenience macros ────────────────────────────────────────────────────────
\newcommand{\Teff}{T_{\!\text{eff}}}
\newcommand{\Eacc}[1]{E_{#1}}
\newcommand{\Prob}[1]{P\!\left(#1\right)}

% ─── Document metadata ─────────────────────────────────────────────────────────
\title{\textbf{Derivation of Expected Accidental Coincidence Rates}\\[6pt]
       \large Muon Shower Detector Array --- Technical Note}
\author{Muon Technical Report v2, \S6 (Expanded Derivation)}
\date{March 2026}

% ═══════════════════════════════════════════════════════════════════════════════
\begin{document}
\maketitle
\tableofcontents
\newpage

% ───────────────────────────────────────────────────────────────────────────────
\section{Overview and Physical Motivation}
% ───────────────────────────────────────────────────────────────────────────────

An \textit{accidental coincidence} is a near-simultaneous signal in
multiple detectors that arises purely from independent, uncorrelated
background events rather than from a genuine cosmic-ray air shower.
Quantifying the expected rate of accidentals is essential: any
statistically significant \emph{excess} of observed coincidences above
this background is evidence for real shower events.

The detector array consists of $N = 8$ scintillator panels (\emph{Nano
detectors}), each generating a time-stamped pulse whenever a charged
particle passes through.  After hardware calibration and a synchronisation
veto to remove periodic artefacts, each detector independently fires at a
roughly constant mean rate.  We ask:

\begin{quote}
\itshape
Given independent Poisson processes, how many times do we expect $k$
specified detectors all to fire within a time window of half-width $\tau$
purely by chance, over an effective run of duration $\Teff$?
\end{quote}

The derivation proceeds in three stages:
\begin{enumerate}
    \item Establish the Poisson framework and define the coincidence window.
    \item Derive the two-fold ($k=2$) accidental rate exactly.
    \item Generalise to an arbitrary $k$-fold coincidence and sum over
          all detector subsets.
\end{enumerate}

% ───────────────────────────────────────────────────────────────────────────────
\section{Definitions and Assumptions}
% ───────────────────────────────────────────────────────────────────────────────

\begin{definition}[Detector rates and run duration]
Let detector $i$ ($i = 1, \ldots, N$, $N=8$) fire according to an
independent, stationary Poisson process with measured rate
$r_i$~(events\,s$^{-1}$).  The effective run duration after dead-time
and veto corrections is $\Teff$~(s).
\end{definition}

\begin{definition}[Coincidence window]
A \emph{coincidence window} of half-width $\tau$ is the symmetric
interval $[t-\tau,\, t+\tau]$ of total width $2\tau$ centred on a
reference time $t$.  Two events are coincident if they fall within the
same window.
\end{definition}

\begin{assumption}[Stationarity]\label{ass:stat}
Each $r_i$ is constant over the run.  This is verified after calibration
and the sync veto.
\end{assumption}

\begin{assumption}[Independence between detectors]\label{ass:indep}
The firing of detector $j$ is statistically independent of the firing of
detector $i \neq j$.  A significant excess of observed coincidences over
the accidental prediction derived below is therefore evidence of
\emph{correlated} events (genuine showers).
\end{assumption}

\begin{assumption}[Rare-event / thin-window approximation]\label{ass:rare}
For the window widths and rates of interest,
\begin{equation}
    r_i \cdot 2\tau \;\ll\; 1.
    \label{eq:rare}
\end{equation}
Concretely, at $r \approx 1$~Hz and $\tau = 0.5$~ms this factor is
$r \cdot 2\tau = 10^{-3}$, accurate to $0.1\,\%$.  Under this condition
the probability of \emph{two or more} events from the same detector
falling inside a single window is $\mathcal{O}\!\left[(r\tau)^2\right]$
and is negligible.
\end{assumption}

% ───────────────────────────────────────────────────────────────────────────────
\section{Poisson Background: Key Probability Results}
% ───────────────────────────────────────────────────────────────────────────────

Before deriving the coincidence rates we collect the Poisson results
that underpin the calculation.

\subsection{Number of events in a fixed interval}

A Poisson process with rate $r$ generates a number of events $X$ in a
time interval of length $\Delta t$ distributed as
\begin{equation}
    \Prob{X = n} = \frac{(r\,\Delta t)^n\, e^{-r\,\Delta t}}{n!},
    \qquad n = 0, 1, 2, \ldots
    \label{eq:poisson-pmf}
\end{equation}
with mean $\langle X \rangle = r\,\Delta t$ and variance $\operatorname{Var}(X)
= r\,\Delta t$.

\subsection{Probability of at least one event}

The probability that at least one event falls in $\Delta t$ is
\begin{equation}
    \Prob{X \geq 1} = 1 - \Prob{X = 0} = 1 - e^{-r\,\Delta t}.
    \label{eq:at-least-one-exact}
\end{equation}
Under the rare-event approximation \eqref{eq:rare} (i.e.\ $r\,\Delta t \ll 1$),
the Taylor expansion of \eqref{eq:at-least-one-exact} gives
\begin{align}
    \Prob{X \geq 1}
        &= 1 - e^{-r\,\Delta t} \notag \\
        &= 1 - \Bigl[1 - r\,\Delta t + \tfrac{(r\,\Delta t)^2}{2} -
           \cdots\Bigr] \notag \\
        &= r\,\Delta t - \tfrac{(r\,\Delta t)^2}{2} + \cdots \notag \\
        &\approx r\,\Delta t.
    \label{eq:at-least-one-approx}
\end{align}
The discarded terms are $\mathcal{O}\!\left[(r\,\Delta t)^2\right]$, which
by Assumption~\ref{ass:rare} are $\lesssim 10^{-6}$ for $r \sim 1$~Hz
and $\tau \leq 10$~ms.

\textbf{Physical interpretation.} In the rare-event regime the window
is so narrow that the probability of getting \emph{any} hit from a
single detector equals simply the mean number of hits expected in the
window.  One can think of $r_j \cdot 2\tau$ as the fractional ``dead
fraction'' of the run that detector $j$'s activity fills.

% ───────────────────────────────────────────────────────────────────────────────
\section{Two-Fold Accidental Coincidence Rate}
\label{sec:two-fold}
% ───────────────────────────────────────────────────────────────────────────────

\subsection{Setup: trigger-and-follower framework}

Consider a specific ordered pair of detectors $(i, j)$.  We use a
\emph{trigger–follower} argument:

\begin{enumerate}
    \item \textbf{Trigger:} Detector $i$ fires at time $t$.  This
          defines a coincidence window $W = [t-\tau,\, t+\tau]$ of
          width $2\tau$.
    \item \textbf{Follower:} We ask whether detector $j$ (independent
          of $i$ by Assumption~\ref{ass:indep}) also fires at least
          once inside $W$.
\end{enumerate}

The key insight is that, conditioned on $i$ having fired, the subsequent
behaviour of $j$ inside $W$ is simply a Poisson process with rate $r_j$
acting over a fixed time interval of length $2\tau$.

\subsection{Probability that the follower fires in the window}

From \eqref{eq:at-least-one-approx}, the probability that detector $j$
fires at least once in the window $W$ is
\begin{equation}
    \Prob{j \text{ fires in } W \mid i \text{ triggered}}
    = 1 - e^{-r_j \cdot 2\tau}
    \;\approx\; r_j \cdot 2\tau,
    \label{eq:pj-in-window}
\end{equation}
where the approximation uses \eqref{eq:at-least-one-approx} and
Assumption~\ref{ass:rare}.

\subsection{Expected number of accidental two-fold coincidences}

Over the entire run, detector $i$ produces a total of
\begin{equation}
    N_i = r_i\,\Teff
    \label{eq:num-triggers}
\end{equation}
trigger events (in expectation).  Each trigger independently opens a
window, and each window has probability \eqref{eq:pj-in-window} of
capturing a hit from $j$.  These trials are:
\begin{itemize}
    \item \textit{Identical:} the same probability applies to every
          window (stationarity, Assumption~\ref{ass:stat}).
    \item \textit{Independent:} non-overlapping triggers are by
          definition independent; even for closely spaced triggers the
          independence follows from the memoryless property of Poisson
          processes.
\end{itemize}
Therefore the expected number of accidental coincidences between the
specific pair $(i, j)$ is
\begin{align}
    \Eacc{2,ij}
        &= \underbrace{r_i\,\Teff}_{\text{no.\ of triggers}}
           \;\times\;
           \underbrace{r_j\cdot 2\tau}_{\text{prob.\ follower hits window}}
           \notag \\[6pt]
        &= r_i\,r_j\,(2\tau)\,\Teff.
        \label{eq:E2ij}
\end{align}

\begin{result}[Two-fold accidental rate, Eq.\ (6.1)]
\begin{equation}
    \boxed{
        \Eacc{2,ij} = r_i\,r_j\,(2\tau)\,\Teff
    }
    \tag{6.1}
\end{equation}
\end{result}

\subsection{Symmetry check}

Expression \eqref{eq:E2ij} is \emph{symmetric} in $i$ and $j$:
swapping which detector is the trigger gives
\[
    \Eacc{2,ji} = r_j\,r_i\,(2\tau)\,\Teff = \Eacc{2,ij}.
\]
This is physically required: the expected number of accidental
coincidences between a pair of detectors cannot depend on our bookkeeping
choice of which one we call the ``trigger.''

\subsection{Double-counting note}

When we later sum over all \emph{pairs}, the combination $\{i,j\}$ is
counted once (as an unordered subset).  We do \textbf{not} add both
$E_{2,ij}$ and $E_{2,ji}$; the trigger is a mathematical device, not a
physical distinction.

% ───────────────────────────────────────────────────────────────────────────────
\section{k-Fold Accidental Coincidence Rate for a Specific Subset}
\label{sec:kfold-specific}
% ───────────────────────────────────────────────────────────────────────────────

\subsection{Generalising the trigger-follower argument}

Let $S = \{i_1, i_2, \ldots, i_k\}$ be a specific ordered set of $k$
detectors.  Designate $i_1$ as the trigger.

\paragraph{Step 1: Trigger events.}
Detector $i_1$ fires $r_{i_1}\Teff$ times (in expectation) over the run.
Each firing at time $t$ opens a window $W = [t-\tau, t+\tau]$.

\paragraph{Step 2: Each follower fires independently.}
For each follower detector $i_m$ ($m = 2, \ldots, k$), the probability
that it fires at least once inside $W$ is (from \eqref{eq:pj-in-window})
\begin{equation}
    \Prob{i_m \text{ fires in } W \mid i_1 \text{ triggered}}
    \approx r_{i_m}\cdot 2\tau.
    \label{eq:follower-prob}
\end{equation}

\paragraph{Step 3: All $k-1$ followers fire simultaneously.}
By Assumption~\ref{ass:indep} (independence between detectors), the
events in \eqref{eq:follower-prob} are independent for different $m$.
The probability that \textit{all} $k-1$ followers fire inside $W$ is
therefore the product of the individual probabilities:
\begin{equation}
    \Prob{\text{all followers fire in }W \mid i_1 \text{ triggered}}
    = \prod_{m=2}^{k} r_{i_m}\cdot 2\tau
    = (2\tau)^{k-1}\prod_{m=2}^{k} r_{i_m}.
    \label{eq:all-followers}
\end{equation}
The $(2\tau)^{k-1}$ factor arises because there are $k-1$ followers each
contributing one power of $2\tau$.

\subsection{Combining trigger count and follower probability}

Multiplying the number of triggers \eqref{eq:num-triggers} by the joint
follower probability \eqref{eq:all-followers}:

\begin{align}
    \Eacc{k,S}
        &= r_{i_1}\Teff
           \;\times\;
           (2\tau)^{k-1}\prod_{m=2}^{k} r_{i_m}
           \notag \\[4pt]
        &= \Teff\,(2\tau)^{k-1}
           \;\underbrace{r_{i_1}\prod_{m=2}^{k} r_{i_m}}_{= \prod_{m=1}^{k} r_{i_m}}
           \notag \\[4pt]
        &= \Teff\,(2\tau)^{k-1}\prod_{m=1}^{k} r_{i_m}.
        \label{eq:Ek-specific}
\end{align}

\begin{result}[$k$-fold accidental rate for a specific $k$-tuple, Eq.\ (6.2)]
\begin{equation}
    \boxed{
        \Eacc{k,S} = \Teff\,(2\tau)^{k-1}\prod_{i \in S} r_i
    }
    \tag{6.2}
\end{equation}
\end{result}

\subsection{Why the choice of trigger drops out}

The final expression \eqref{eq:Ek-specific} contains the product of
\emph{all} $k$ rates $r_{i_1}\cdots r_{i_k}$, regardless of which
detector was labelled the trigger.  The trigger contributed one factor of
$r_{i_1}$ to the count $N_{i_1} = r_{i_1}\Teff$, while the $k-1$
followers each contributed one factor of $r_{i_m}\cdot 2\tau$.  The
product simply collects all $k$ rates symmetrically.  In other words, had
we chosen $i_2$ as the trigger instead, the derivation would have given
\[
    r_{i_2}\Teff \times (2\tau)^{k-1}\prod_{\substack{m=1\\m\neq 2}}^{k} r_{i_m}
    = \Teff\,(2\tau)^{k-1}\prod_{m=1}^{k} r_{i_m},
\]
which is identical.  The result is fully symmetric in all $k$ detectors.

\subsection{Dimensional analysis}

\begin{center}
\renewcommand{\arraystretch}{1.3}
\begin{tabular}{lll}
\toprule
Quantity & Units & Exponent in \eqref{eq:Ek-specific} \\
\midrule
$\Teff$      & s                    & 1 \\
$(2\tau)^{k-1}$ & s$^{k-1}$         & $k-1$ \\
$\prod r_{i_m}$ & (s$^{-1}$)$^k$  & $k$ \\
\midrule
Product      & s$^{1+(k-1)-k}$ = s$^0$ = dimensionless & \\
\bottomrule
\end{tabular}
\end{center}

The result is a dimensionless \emph{count}, as required.

% ───────────────────────────────────────────────────────────────────────────────
\section{Total Expected k-Fold Accidentals (All Subsets)}
\label{sec:total-kfold}
% ───────────────────────────────────────────────────────────────────────────────

\subsection{Summing over all k-element subsets}

The array contains $N = 8$ detectors.  The number of distinct
$k$-element subsets (unordered combinations) is $\binom{N}{k} = \binom{8}{k}$.
The total expected number of accidental $k$-fold coincidences is the
sum of \eqref{eq:Ek-specific} over all such subsets:
\begin{equation}
    \Eacc{k}
    = \sum_{\substack{S\,\subseteq\,\{1,\ldots,8\}\\ |S|=k}}
        \Teff\,(2\tau)^{k-1}\prod_{i\in S} r_i.
    \label{eq:Ek-total-pre}
\end{equation}
Since $\Teff$ and $(2\tau)^{k-1}$ do not depend on the subset $S$, they
factor out of the sum:
\begin{equation}
    \boxed{
        \Eacc{k}
        = \Teff\,(2\tau)^{k-1}
          \sum_{\substack{S\,\subseteq\,\{1,\ldots,8\}\\ |S|=k}}
          \prod_{i\in S} r_i.
    }
    \tag{6.3}
    \label{eq:Ek-total}
\end{equation}

\begin{result}[Total $k$-fold accidental rate, Eq.\ (6.3)]
Equation~\eqref{eq:Ek-total} is exact within the Poisson framework and
the rare-event approximation.  It is implemented numerically in
\texttt{expected\_accidentals()} using the elementary symmetric
polynomial $e_k(r_1,\ldots,r_8)$ to enumerate all $\binom{8}{k}$ subsets
efficiently.
\end{result}

\subsection{Connection to elementary symmetric polynomials}

The sum $\displaystyle\sum_{|S|=k}\prod_{i\in S} r_i$ is precisely the
$k$-th \emph{elementary symmetric polynomial} in the rates:
\[
    e_k(r_1,\ldots,r_N) = \sum_{1\le i_1 < i_2 < \cdots < i_k \le N}
        r_{i_1} r_{i_2} \cdots r_{i_k}.
\]
For instance,
\[
    e_2(r_1,\ldots,r_8) = \sum_{i<j} r_i r_j, \qquad
    e_3(r_1,\ldots,r_8) = \sum_{i<j<\ell} r_i r_j r_\ell, \quad\ldots
\]
so that $\Eacc{k} = \Teff(2\tau)^{k-1}\,e_k(r_1,\ldots,r_8)$.

\subsection{Scaling with fold order}

A key insight is how $\Eacc{k}$ falls as $k$ increases.  For
simplicity suppose all rates are equal: $r_i = r$.  Then
$e_k(r,\ldots,r) = \binom{8}{k}r^k$, so
\[
    \Eacc{k} = \Teff\,(2\tau)^{k-1}\,\binom{8}{k}\,r^k
    = \Teff\,r\,\binom{8}{k}\,(r\cdot 2\tau)^{k-1}.
\]
The ratio of consecutive fold orders is
\begin{equation}
    \frac{\Eacc{k+1}}{\Eacc{k}}
    = \frac{\binom{8}{k+1}}{\binom{8}{k}}\cdot r\cdot 2\tau
    = \frac{8-k}{k+1}\cdot r\cdot 2\tau.
    \label{eq:fold-ratio}
\end{equation}
At $r = 1$~Hz and $2\tau = 1$~ms, the factor $r\cdot 2\tau = 10^{-3}$,
so each additional fold suppresses $\Eacc{k}$ by roughly three orders of
magnitude (multiplied by a combinatorial factor of order unity).  This
is why a single observed 5-fold event, when $\Eacc{5} \sim 10^{-3}$ at
the 1~ms window, constitutes an enormous statistical excess.

% ───────────────────────────────────────────────────────────────────────────────
\section{Numerical Results for the February 27 Dataset}
\label{sec:numerical}
% ───────────────────────────────────────────────────────────────────────────────

The following table reproduces the expected accidentals from the
Technical Report (\S6.3) for the February 27, 2026 run
($\Teff \approx 3{,}576$~s, rates $r_i \in [0.81, 1.25]$~Hz):

\begin{table}[h!]
\centering
\caption{Expected accidental coincidence counts $\Eacc{k}$ from
Eq.~\eqref{eq:Ek-total} for selected coincidence windows and fold orders.
Values are rounded to two significant figures.}
\label{tab:accidentals}
\renewcommand{\arraystretch}{1.4}
\begin{tabular}{>{\centering\arraybackslash}p{1.5cm}
                >{\centering\arraybackslash}p{1.8cm}
                >{\centering\arraybackslash}p{1.8cm}
                >{\centering\arraybackslash}p{1.8cm}
                >{\centering\arraybackslash}p{1.8cm}
                >{\centering\arraybackslash}p{1.8cm}}
\toprule
$2\tau$ & $E_2$ & $E_3$ & $E_4$ & $E_5$ & $E_6$ \\
\midrule
1~ms  & 202    & 0.80   & $<10^{-2}$  & $<10^{-3}$  & $<10^{-4}$ \\
2~ms  & 405    & 3.2    & $<10^{-2}$  & $<10^{-3}$  & $<10^{-4}$ \\
5~ms  & 1\,012  & 20.3   & 0.3         & $<10^{-2}$  & $<10^{-3}$ \\
10~ms & 2\,024  & 81.1   & 2.0         & 0.02        & $<10^{-3}$ \\
\bottomrule
\end{tabular}
\end{table}

To illustrate the scaling law \eqref{eq:fold-ratio} concretely, at
$2\tau = 1$~ms:
\[
    \frac{\Eacc{3}}{\Eacc{2}}
    \approx \frac{0.80}{202} \approx 4\times 10^{-3}
    \approx \frac{7}{2}\cdot r_{\text{avg}}\cdot 2\tau
    = \frac{7}{2}\cdot 1\,\text{Hz}\cdot 10^{-3}\,\text{s},
\]
consistent with \eqref{eq:fold-ratio} at $k=2$.

% ───────────────────────────────────────────────────────────────────────────────
\section{Validity Conditions}
\label{sec:validity}
% ───────────────────────────────────────────────────────────────────────────────

The derivation above rests on three conditions that must hold in practice.

\paragraph{1.~Poisson statistics (stationarity and independence within a detector).}
Each detector must fire at a constant mean rate with no clustering or
periodicity in its own pulse stream.  The synchronisation veto removes
hardware-generated periodic artefacts, making the post-veto event stream
consistent with Poisson statistics.  Stationarity is verified by checking
that the measured rate $r_i$ is stable across short sub-runs.

\paragraph{2.~Rare-event / thin-window approximation.}
Equation \eqref{eq:rare} must hold: $r_i \cdot 2\tau \ll 1$.  For the
rates and windows in Table~\ref{tab:accidentals}:

\begin{center}
\renewcommand{\arraystretch}{1.3}
\begin{tabular}{ccc}
\toprule
$2\tau$ & $r_{\max}\cdot 2\tau$ & Fractional error in $\Eacc{k}$ \\
\midrule
1~ms  & $1.25\times10^{-3}$ & $\sim 0.1\,\%$ \\
10~ms & $1.25\times10^{-2}$ & $\sim 1\,\%$ \\
\bottomrule
\end{tabular}
\end{center}

At all windows used, the approximation error is well within the
statistical uncertainty of the measurements.

\paragraph{3.~Independence between detectors.}
Equation~\eqref{eq:all-followers} multiplies probabilities for different
detectors using the product rule, which requires independence
(Assumption~\ref{ass:indep}).  Physically this holds if the detectors do
not share a noise source and if cosmic-ray showers — which would create
correlated hits — are rare.  The entire point of the accidental
calculation is to define a baseline against which genuine shower
correlations are detected as a statistically significant excess.

% ───────────────────────────────────────────────────────────────────────────────
\section{Summary of Results}
\label{sec:summary}
% ───────────────────────────────────────────────────────────────────────────────

\begin{result}[Two-fold accidentals between a specific pair $(i,j)$]
\[
    \Eacc{2,ij} = r_i\,r_j\,(2\tau)\,\Teff. \tag{6.1}
\]
\end{result}

\begin{result}[$k$-fold accidentals for a specific $k$-tuple $S$]
\[
    \Eacc{k,S} = \Teff\,(2\tau)^{k-1}\prod_{i\in S} r_i. \tag{6.2}
\]
\end{result}

\begin{result}[Total $k$-fold accidentals summed over all subsets]
\[
    \Eacc{k}
    = \Teff\,(2\tau)^{k-1}
      \sum_{\substack{S\,\subseteq\,\{1,\ldots,8\}\\ |S|=k}}
      \prod_{i\in S} r_i
    = \Teff\,(2\tau)^{k-1}\,e_k(r_1,\ldots,r_8).
    \tag{6.3}
\]
The steep suppression with $k$ — approximately a factor of $r\cdot 2\tau
\approx 10^{-3}$ per additional fold at 1~ms windows — means that even a
single observed high-fold coincidence event constitutes an overwhelming
statistical excess above this accidental background.
\end{result}

\end{document}
